% archon:covers AlgebraicJacobian/RiemannRoch/Adelic/Substrate.lean
% archon:covers AlgebraicJacobian/RiemannRoch/Adelic/Cokernel.lean
% archon:covers AlgebraicJacobian/RiemannRoch/Adelic/FinitenessP1.lean
% archon:covers AlgebraicJacobian/RiemannRoch/Adelic/ChiLedger.lean
% archon:covers AlgebraicJacobian/RiemannRoch/Adelic/Duality.lean
\chapter{Adelic Riemann--Roch: the repartition cokernel (RR.A)}
\label{chap:RiemannRoch_Adelic}

\section{Provenance and standing conventions}
\label{sec:RiemannRoch_Adelic_intro}

This chapter builds the \emph{adelic} (Weil-repartition) approach to
Riemann--Roch for the smooth proper geometrically integral curve
\(C / k\) of the challenge, over an \emph{arbitrary} base field \(k\)
(not assumed algebraically closed). It reuses the divisor substrate of
\Cref{chap:RiemannRoch_WeilDivisor} (\(\mathrm{PrimeDivisor}\), the
valuation \(\mathrm{order}\), the degree map, the class
\(\mathrm{IsRegularInCodimensionOne}\)) and the sheaf-cohomology
substrate of \Cref{chap:Cohomology_MayerVietoris} and
\Cref{chap:Genus} (the structure sheaf as a sheaf of \(k\)-modules
\(\toModuleKSheaf C\), its cohomology \(\HModule\), the Mayer--Vietoris
long exact sequence, and the Čech carriers
\(\mathrm{IsCechAcyclicCover}\) / \(\HasCechToHModuleIso\)).

\paragraph{The concrete-repartition idea.}
Classically (Weil), for the function field \(K = k(C)\) one forms the
\emph{adele} (repartition) space
\(\mathbb{A}_K = \{(x_P)_P \in \prod_P K : v_P(x_P) \ge 0 \text{ for
almost all } P\}\), its bounded subspaces
\(\mathbb{A}_K(D) = \{(x_P) : v_P(x_P) \ge -D(P)\ \forall P\}\), and the
two Riemann--Roch spaces
\[
  L(D) = K \cap \mathbb{A}_K(D)
        = \{ f \in K : \operatorname{div} f + D \ge 0\},
  \qquad
  H^1(D) = \mathbb{A}_K / (\mathbb{A}_K(D) + K).
\]
The genus is \(g = \dim_k H^1(0)\); the index of specialty is
\(i(D) = \dim_k H^1(D)\). We do \emph{not} build \(\mathbb{A}_K\): its
Lean assembly is absent from mathlib in the function-field case, and the
comparison of \(\mathbb{A}_K/K\) with sheaf cohomology re-introduces a
surjectivity/existence gap. Instead we use the observation that for a
\emph{fixed \(2\)-element affine cover} the ``adele-for-that-cover''
first cohomology \emph{is} the Čech \(H^1\) definitionally --- a single
cokernel, no colimit and no existence theorem.

\paragraph{The \(2\)-affine-cover cokernel.}
Fix an affine Mayer--Vietoris square
\(S : C.\mathrm{left}.\AffineCoverMVSquare\) (\Cref{chap:Cohomology_MayerVietoris}):
opens \(U_0, U_1 \subseteq C\) with \(U_0, U_1, U_0 \cap U_1\) all affine
and \(U_0 \sqcup U_1 = \top\). Write, as \(k\)-subspaces of \(K\) via the
fraction-field embedding,
\[
  A_i = \Gamma(U_i, \struct{C}), \qquad
  A_{01} = \Gamma(U_0 \cap U_1, \struct{C}).
\]
The degree-\(1\) Čech cohomology of the two-element cover is the cokernel
of the difference map \((f_0, f_1) \mapsto f_0|_{01} - f_1|_{01}\):
\[
  \check{H}^1 \;=\; A_{01} \big/ (A_0 + A_1).
\]
For a Weil divisor \(D\) the \emph{twisted} version replaces each ring of
sections by the Riemann--Roch sheaf of sections
\(\Gamma(U, \struct{C}(D)) = \{ f \in K : v_P(f) \ge -D(P) \text{ for all
prime divisors } P \in U\}\), giving
\[
  L(D) = \Gamma(U_0, \struct{C}(D)) \cap \Gamma(U_1, \struct{C}(D)),
  \qquad
  \check{H}^1(D) = \Gamma(U_0 \cap U_1, \struct{C}(D)) \big/
    \bigl(\Gamma(U_0, \struct{C}(D)) + \Gamma(U_1, \struct{C}(D))\bigr).
\]
These are the algebraic-geometry incarnations of Weil's \(L(D)\) and
\(H^1(D)\), and the bridge of \Cref{sec:RiemannRoch_Adelic_cokernel}
identifies \(\check{H}^1(0)\) with the project's
\(\HModule\, k\, (\toModuleKSheaf C)\, 1 = H^1(C, \struct{C})\).

\paragraph{Residue-degree bookkeeping.}
Because \(k\) need not be algebraically closed, every place carries a
\emph{residue degree}. For a prime divisor \(P\) with residue field
\(\kappa(P) = \struct{C, P}/\mathfrak{m}_P\), set
\[
  \deg P \;=\; [\kappa(P) : k] \;=\; \dim_k \kappa(P) \;<\; \infty,
  \qquad
  \deg D \;=\; \sum_P D(P)\,\deg P .
\]
Over an algebraically closed \(k\) one has \(\deg P = 1\) and this
recovers the classical count; the whole \(\chi\)-ledger below carries the
residue degrees. This is the degree of
\Cref{def:divisor_degree,thm:divisor_degree_hom}.

\paragraph{Notation.} We write \(\ell(D) = \dim_k L(D)\),
\(i(D) = \dim_k \check{H}^1(D)\), and
\(\chi(D) = \ell(D) - i(D)\) for the Euler characteristic. The genus is
\(g = i(0) = \dim_k \check{H}^1(0)\) (\Cref{def:genus}, once
\Cref{thm:adelic_module_finite_genus} makes it a finite number).

\paragraph{Tiering.} The development is layered so that the earlier tiers
verify without the later ones.
\emph{Tier 0} (\Cref{sec:RiemannRoch_Adelic_substrate}) records the
Dedekind chart, the point valuation, the sheaf of sections, and the
finite-support lemma.
\emph{Tier 1} (\Cref{sec:RiemannRoch_Adelic_cokernel,sec:RiemannRoch_Adelic_finiteness})
delivers the single live-tree payoff --- finite-dimensionality of
\(H^1(C, \struct{C})\), hence an honest \(\genus C\) --- funnelling
through the keystone
\(\Cref{thm:adelic_finiteDimensional_cokernel_zero}\).
\emph{Tier 2} (\Cref{sec:RiemannRoch_Adelic_chi}) is the
\(\chi\)-ledger: the Riemann inequality and finiteness of every
\(\check{H}^1(D)\), keystone
\(\Cref{thm:adelic_chi_eq_chi_zero_add_degree}\).
\emph{Tier 3} (\Cref{sec:RiemannRoch_Adelic_duality}) is residue-pairing
Serre duality and the vanishing \(i(D) = 0\) for \(\deg D > 2g - 2\); its
two deepest nodes are gated by single-field \(\mathrm{Prop}\)-classes, in
the style of \(\mathrm{HasPicScheme}\), so the earlier tiers are
unaffected.

\medskip
Throughout, \(C : \mathrm{Over}\,(\Spec (.of\ k))\) with
\([\mathrm{IsProper}\ C.\hom]\),
\([\mathrm{SmoothOfRelativeDimension}\ 1\ C.\hom]\) and
\([\mathrm{GeometricallyIrreducible}\ C.\hom]\), matching the standing
hypotheses of \Cref{def:genus}; \(K = C.\mathrm{left}.\mathrm{functionField}\).

\section{Tier 0 --- the adelic substrate}
\label{sec:RiemannRoch_Adelic_substrate}

\begin{proposition}
[Affine chart is a Dedekind domain with fraction field \(K\)]
  \label{thm:adelic_chartRing_isDedekindDomain}
  \lean{AlgebraicGeometry.Adelic.chartRing_isDedekindDomain}
  \uses{def:isRegularInCodimensionOne}
  % [scheidler] valuation rings / DVR foundations; [cff] Prop. 4.2 (DVR characterisation)
  \source{scheidler-function-fields, cff-curves-function-fields}
  Let \(U \subseteq C\) be a nonempty affine open with coordinate ring
  \(A_U = \Gamma(U, \struct{C})\). Then \(A_U\) is a Dedekind domain and
  the canonical map \(A_U \hookrightarrow K\) exhibits \(K\) as the field
  of fractions of \(A_U\).
\end{proposition}
\begin{proof}
  Since \(C\) is integral, \(A_U\) is an integral domain, and for a
  nonempty affine open of an integral scheme the natural map to the
  function field \(A_U \to K\) is injective with
  \(\mathrm{IsFractionRing}\, A_U\, K\) (the localisation of \(A_U\) at
  its generic point is \(K\)); this is the affine-chart fraction-field
  statement used to build \(\texttt{functionFieldIso}\)
  (\Cref{def:functionFieldIso}). It remains to verify the Dedekind
  axioms. \(A_U\) is Noetherian: \(C\) is of finite type over the field
  \(k\), hence locally Noetherian, so \(A_U\) is a finitely generated
  \(k\)-algebra and Noetherian. \(A_U\) is a domain that is not a field
  (a curve has closed points), of Krull dimension one: every nonzero
  prime of \(A_U\) is maximal because the corresponding point of \(U\)
  has codimension one in the one-dimensional \(C\). At each such
  height-one prime \(\mathfrak p\), the localisation
  \((A_U)_{\mathfrak p} = \struct{C, x}\) is the stalk at the
  corresponding codimension-one point \(x\), which is a discrete
  valuation ring by \(\mathrm{IsRegularInCodimensionOne}\)
  (\Cref{def:isRegularInCodimensionOne}), automatic here from the
  smoothness of \(C\). A one-dimensional Noetherian domain whose
  localisations at height-one primes are discrete valuation rings is
  integrally closed, hence Dedekind by the standard characterisation
  (a Noetherian integrally closed domain of dimension \(\le 1\)).
  \paragraph{Gate.} The one non-formal input is ``smooth curve chart
  \(\Rightarrow\) regular chart'', supplied here by
  \(\mathrm{IsRegularInCodimensionOne}\). Should the passage from
  smoothness of \(C.\hom\) to that class be unavailable at a use site, it
  is threaded as the single-field class
  \(\mathrm{HasDedekindChart}\, C\), discharged by \(\texttt{inferInstance}\)
  once the regularity instance is present.
\end{proof}

\begin{definition}
\leanok
[Point valuation and residue degree]
  \label{def:adelic_pointValuation}
  \lean{AlgebraicGeometry.Adelic.pointValuation}
  \uses{thm:adelic_chartRing_isDedekindDomain, def:order_at_point, lem:order_eq_order_restrict}
  % [scheidler] p. 24 (deg P = [F_P : K]); [cff] (C2) deg P = dim_K O_P/P
  \source{scheidler-function-fields, cff-curves-function-fields}
  Let \(P\) be a prime divisor whose point lies in a nonempty affine open
  \(U\), with Dedekind chart \(A_U\)
  (\Cref{thm:adelic_chartRing_isDedekindDomain}). The prime \(P\)
  determines a height-one prime \(\mathfrak p_P\) of \(A_U\), and the
  discrete valuation of the challenge,
  \(v_P = \mathrm{order}\, P\) on \(K\) (\Cref{def:order_at_point}),
  coincides with the \(\mathfrak p_P\)-adic valuation of \(A_U\) on
  \(K = \operatorname{Frac} A_U\). Its residue field is
  \(\kappa(P) = A_U/\mathfrak p_P = \struct{C, P}/\mathfrak m_P\), and the
  \emph{degree} of \(P\) is the finite number
  \(\deg P = \dim_k \kappa(P) = [\kappa(P) : k]\).
\end{definition}
\begin{proof}
  The stalk \(\struct{C, P}\) is the localisation \((A_U)_{\mathfrak p_P}\)
  because \(P\) lies in the affine \(U\); its normalised valuation, whose
  \(\mathrm{WithZero.log}\) is \(\mathrm{order}\, P\)
  (\Cref{def:order_at_point}), is by definition the
  \(\mathfrak p_P\)-adic valuation. Independence of the chosen chart is
  the naturality \Cref{lem:order_eq_order_restrict}: \(\mathrm{order}\)
  is invariant under open-immersion restriction, so any two affine charts
  containing \(P\) give the same valuation on \(K\). Finiteness of
  \(\deg P\): \(\kappa(P)\) is the residue field of a closed point of the
  finite-type \(k\)-scheme \(C\); choosing a nonconstant \(x \in K\)
  bounds \([\kappa(P) : k] \le [K : k(x)] < \infty\).
\end{proof}

\begin{definition}
\leanok
[Sheaf of sections of a divisor; the Riemann--Roch space]
  \label{def:adelic_sectionOfDivisor}
  \lean{AlgebraicGeometry.Adelic.sectionOfDivisor}
  \uses{def:adelic_pointValuation, def:codim1_cycles, lem:order_mul_of_ne_zero}
  % [papaioannou] Def. 1.9 (L(D)); [cff] sec. 12 (L(A), Riemann-Roch space)
  \source{papaioannou-algebraic-rr, cff-curves-function-fields}
  For an open \(U \subseteq C\) and a Weil divisor \(D\)
  (\Cref{def:codim1_cycles}) define the additive subgroup of \(K\)
  \[
    \Gamma(U, \struct{C}(D)) \;=\;
      \{ f \in K : f = 0 \text{ or } v_P(f) \ge -D(P) \text{ for every
         prime divisor } P \text{ with point in } U\}
  \]
  (the disjunct \(f = 0\) encodes \(v_P(0) = +\infty\)). It is in fact a
  \(k\)-subspace, since scaling by \(c \in k^\times\) preserves every
  order \(v_P(cf) = v_P(f)\). Its global instance for the fixed cover
  gives the \emph{Riemann--Roch space}
  \(L(D) = \Gamma(U_0, \struct{C}(D)) \cap \Gamma(U_1, \struct{C}(D))\),
  equal to \(\{ f \in K : \operatorname{div} f + D \ge 0\}\) since
  \(U_0 \cup U_1 = C\).
\end{definition}
\begin{proof}
  The set is a \(k\)-submodule of \(K\). It contains \(0\)
  (\(v_P(0) = +\infty \ge -D(P)\)). It is closed under addition by the
  ultrametric inequality of the discrete valuation
  \(v_P(f + h) \ge \min(v_P f, v_P h) \ge -D(P)\). It is closed under
  scaling by \(c \in k^\times\): a nonzero constant is a unit in every
  stalk, so \(v_P(c) = 0\) and \(v_P(cf) = v_P(c) + v_P(f) = v_P(f)\) by
  \(\Cref{lem:order_mul_of_ne_zero}\) (with \(v_P(0\cdot f) = +\infty\)).
  The intersection identity for \(L(D)\) holds because a prime divisor of
  \(C\) has its point in \(U_0\) or in \(U_1\); imposing \(v_P(f) \ge
  -D(P)\) on the union of the two chart conditions is imposing it at
  every prime divisor, i.e.\ \(\operatorname{div} f + D \ge 0\).
\end{proof}

\begin{lemma}
[Finiteness of the non-integral places on a Dedekind chart]
  \label{thm:adelic_finite_support_order}
  \lean{AlgebraicGeometry.Adelic.finite_places_ne_one}
  \leanok
  \uses{thm:adelic_chartRing_isDedekindDomain}
  % [scheidler] p. 31 (finiteness of zeros/poles); [papaioannou] Cor. 1.12
  \source{scheidler-function-fields, papaioannou-algebraic-rr}
  Let \(R\) be a Dedekind domain with fraction field \(K\). For a nonzero
  \(k \in K\) the set \(\{ v \in \operatorname{Spec}^1 R : v(k) \ne 1\}\) of
  height-one places at which \(k\) is a zero or a pole is finite.

  This is the affine-chart core of the finite-support statement: a curve is
  covered by finitely many Dedekind charts, and applying the lemma on each
  chart shows that for \(f \in K^\times\) only finitely many prime divisors
  have \(v_P(f) \ne 0\), so \(\operatorname{div} f = \sum_P v_P(f)\,P\) is a
  well-defined Weil divisor. The general scheme-level assembly (finite affine
  cover plus the height-one/minimal-primes bound) is
  \Cref{lem:rationalMap_order_finite_support}.
\end{lemma}
\begin{proof}
  \leanok
  Apply the finiteness of the support of a nonzero fractional ideal in a
  Dedekind domain to both \(k\) and \(k^{-1}\). A place \(v\) with
  \(v(k) \ne 1\) satisfies either \(v(k) > 1\) (a zero of \(k\)) or
  \(v(k) < 1\), equivalently \(v(k^{-1}) > 1\) (a pole of \(k\)); each of
  the two sets of such places is finite, and their union contains every
  place with \(v(k) \ne 1\).
\end{proof}

\section{Tier 1a --- the cover cokernel and the bridge to \(H^1\)}
\label{sec:RiemannRoch_Adelic_cokernel}

\begin{definition}
[The twisted cover cohomology]
  \label{def:adelic_coverH1}
  \lean{AlgebraicGeometry.Adelic.coverH1}
  \uses{def:adelic_sectionOfDivisor, def:Scheme_AffineCoverMVSquare}
  For a Weil divisor \(D\) and the fixed affine square \(S\), the
  \emph{cover cohomology} is the cokernel of the difference map
  \[
    d_D \colon
      \Gamma(U_0, \struct{C}(D)) \oplus \Gamma(U_1, \struct{C}(D))
      \longrightarrow \Gamma(U_0 \cap U_1, \struct{C}(D)),
    \qquad (f_0, f_1) \mapsto f_0 - f_1,
  \]
  (restrictions to \(U_0 \cap U_1\), all viewed inside \(K\)):
  \[
    \check{H}^1(D) \;=\;
      \Gamma(U_0 \cap U_1, \struct{C}(D)) \big/
      \bigl(\Gamma(U_0, \struct{C}(D)) + \Gamma(U_1, \struct{C}(D))\bigr).
  \]
  Its kernel is \(\ker d_D = L(D)\) (\Cref{def:adelic_sectionOfDivisor}).
  We write \(\check{H}^1 = \check{H}^1(0)\), the untwisted cokernel
  \(A_{01}/(A_0 + A_1)\).
\end{definition}
\begin{proof}
  The two summands \(\Gamma(U_i, \struct{C}(D))\) and the target
  \(\Gamma(U_0 \cap U_1, \struct{C}(D))\) are \(k\)-submodules of \(K\)
  (\Cref{def:adelic_sectionOfDivisor}), and \(d_D\) is \(k\)-linear;
  hence the image \(A_0(D) + A_1(D)\) is a submodule and the quotient is
  a \(k\)-module. An element \((f_0, f_1)\) lies in \(\ker d_D\) iff
  \(f_0 = f_1\) in \(K\) with \(f_0 \in \Gamma(U_0, \struct{C}(D))\) and
  \(f_1 \in \Gamma(U_1, \struct{C}(D))\), i.e.\ iff the common value lies
  in \(\Gamma(U_0, \struct{C}(D)) \cap \Gamma(U_1, \struct{C}(D)) =
  L(D)\).
\end{proof}

\begin{proposition}
[Two-element Čech \(H^1\) is the cokernel]
  \label{thm:adelic_cechH1_eq_cokernel}
  \lean{AlgebraicGeometry.Adelic.cechH1_eq_cokernel}
  \uses{def:adelic_coverH1, def:Scheme_cechCohomology, def:Scheme_toModuleKSheaf}
  There is a \(k\)-linear isomorphism
  \[
    \mathrm{cechCohomology}\, C\, (\toModuleKSheaf C)\, S.\cover\, 1
    \;\simeq_k\; \check{H}^1 \;=\; A_{01}/(A_0 + A_1).
  \]
\end{proposition}
\begin{proof}
  For the two-element cover \(\{U_0, U_1\}\) the (alternating) Čech
  cochain complex of \(\toModuleKSheaf C\) is concentrated in degrees
  \(0\) and \(1\): the \(0\)-cochains are
  \(C^0 = \Gamma(U_0) \oplus \Gamma(U_1) = A_0 \oplus A_1\), the
  \(1\)-cochains are \(C^1 = \Gamma(U_0 \cap U_1) = A_{01}\), and
  \(C^{\ge 2} = 0\) since every triple index repeats a member of a
  \(2\)-element set. The differential \(C^0 \to C^1\) is
  \((f_0, f_1) \mapsto f_0|_{01} - f_1|_{01}\), so
  \(H^1 = C^1/\operatorname{im} = A_{01}/(A_0 + A_1)\); the image is
  \(A_0 + A_1\) because each \(A_i\) is a \(k\)-subspace closed under
  negation, so
  \(\{ f_0 - f_1 : f_i \in A_i\} = A_0 + A_1\) inside \(A_{01}\).
\end{proof}

\begin{proposition}
[\(\mathrm{IsCechAcyclicCover}\) for the affine square]
  \label{thm:adelic_instIsCechAcyclicCover}
  \lean{AlgebraicGeometry.Adelic.instIsCechAcyclicCover}
  \uses{def:Scheme_IsCechAcyclicCover, lem:affine_serre_vanishing, lem:affine_injective_acyclic}
  The two-element affine cover \(S.\cover = \{U_0, U_1\}\) satisfies
  \(\mathrm{IsCechAcyclicCover}\, (\toModuleKSheaf C)\, S.\cover\): the
  higher Čech cohomology of the structure sheaf along the cover
  vanishes.
\end{proposition}
\begin{proof}
  Acyclicity of an affine cover for a quasi-coherent sheaf is the
  Serre-vanishing seed of \Cref{lem:affine_serre_vanishing}: on an affine
  open the higher cohomology of \(\toModuleKSheaf C\) vanishes, and the
  injective-acyclic reduction \Cref{lem:affine_injective_acyclic}
  propagates this to the Čech complex of the cover whose members and all
  of whose finite intersections are affine. For the two-element square
  the only nonempty intersections are \(U_0, U_1, U_0 \cap U_1\), all
  affine by construction of \(\AffineCoverMVSquare\)
  (\Cref{def:Scheme_AffineCoverMVSquare}); the vanishing seed applies to
  each, giving \(\mathrm{IsCechAcyclicCover}\).
\end{proof}

\begin{proposition}
[Čech-to-derived comparison for the affine square]
  \label{thm:adelic_instHasCechToHModuleIso}
  \lean{AlgebraicGeometry.Adelic.instHasCechToHModuleIso}
  \uses{def:Scheme_HasCechToHModuleIso, thm:adelic_instIsCechAcyclicCover, thm:Scheme_AffineCoverMVSquare_HModule_prime_sequence_exact}
  The affine square satisfies
  \(\HasCechToHModuleIso\, (\toModuleKSheaf C)\, S.\cover\): the Čech
  cohomology of the cover computes the sheaf cohomology of
  \(\toModuleKSheaf C\).
\end{proposition}
\begin{proof}
  The comparison is furnished by the Mayer--Vietoris presentation of the
  affine square. The exact sequence
  \Cref{thm:Scheme_AffineCoverMVSquare_HModule_prime_sequence_exact}
  expresses \(\HModulePrime\) at the join \(U_0 \sqcup U_1 = \top\) as the
  cohomology of the two-term complex
  \(\HModulePrime(U_0) \oplus \HModulePrime(U_1) \to
  \HModulePrime(U_0 \cap U_1)\); since the cover is Čech-acyclic
  (\Cref{thm:adelic_instIsCechAcyclicCover}) the higher terms
  \(H^{>0}(U_i)\) vanish, so this two-term complex is exactly the Čech
  complex, and its cohomology agrees with the derived-functor cohomology
  \(\HModulePrime\) degreewise. Packaging the degreewise isomorphisms is
  \(\HasCechToHModuleIso\).
  \paragraph{Gate.} This is the substantive Čech-to-derived comparison.
  Where its Lean discharge is not yet available it is threaded as the
  existing carrier class \(\HasCechToHModuleIso\), supplied at the use
  site and later closed by \(\texttt{inferInstance}\); every downstream consumer
  verifies against the hypothesis meanwhile.
\end{proof}

\begin{theorem}
[Bridge: \(H^1(C, \struct{C})\) is the cover cokernel]
  \label{thm:adelic_HModuleOne_linearEquiv_cokernel}
  \lean{AlgebraicGeometry.Adelic.HModuleOne_linearEquiv_cokernel}
  \uses{thm:adelic_cechH1_eq_cokernel, thm:adelic_instIsCechAcyclicCover, thm:adelic_instHasCechToHModuleIso, def:Scheme_HModule_prime_eq_HModule_linearEquiv, def:Scheme_cechToHModuleIso, def:Scheme_HModule}
  There is a \(k\)-linear isomorphism
  \[
    \HModule\, k\, (\toModuleKSheaf C)\, 1
    \;\simeq_k\; \check{H}^1 \;=\; A_{01}/(A_0 + A_1).
  \]
\end{theorem}
\begin{proof}
  Chain three isomorphisms.
  \begin{enumerate}
    \item The universe bridge
      \Cref{def:Scheme_HModule_prime_eq_HModule_linearEquiv} identifies
      \(\HModule\, k\, (\toModuleKSheaf C)\, 1\) with
      \(\HModulePrime\, k\, (\toModuleKSheaf C)\, 1\) at the top open,
      moving between the two universe levels at which the project's
      cohomology is defined.
    \item Under the discharged carriers
      \(\mathrm{IsCechAcyclicCover}\)
      (\Cref{thm:adelic_instIsCechAcyclicCover}) and
      \(\HasCechToHModuleIso\)
      (\Cref{thm:adelic_instHasCechToHModuleIso}), the comparison
      \(\texttt{cechToHModuleIso}\) (\Cref{def:Scheme_cechToHModuleIso})
      identifies \(\HModulePrime\, 1\) at
      \(\bigsqcup S.\cover = \top\) with
      \(\mathrm{cechCohomology}\, C\, (\toModuleKSheaf C)\, S.\cover\, 1\).
    \item The two-element identification
      \Cref{thm:adelic_cechH1_eq_cokernel} rewrites the Čech
      \(H^1\) as the cokernel \(A_{01}/(A_0 + A_1)\).
  \end{enumerate}
  Composing the three \(k\)-linear isomorphisms gives the claim.
\end{proof}

\section{Tier 1b --- finiteness of the base cokernel via \(\mathbb{P}^1\)}
\label{sec:RiemannRoch_Adelic_finiteness}

The bridge \Cref{thm:adelic_HModuleOne_linearEquiv_cokernel} reduces
finite-dimensionality of \(H^1(C, \struct{C})\) to that of the concrete
cokernel \(\check{H}^1 = A_{01}/(A_0 + A_1)\). Unlike the
differential-geometry sibling, which obtained base-case finiteness from
compactness (finite-dimensional Dolbeault cohomology), we must prove it
by hand. The route is reduction to \(\mathbb{P}^1\) through a finite
morphism.

\begin{proposition}
[A finite map to \(\mathbb{P}^1\)]
  \label{thm:adelic_exists_finiteMorphismToP1}
  \notready  % gate class AlgebraicGeometry.Adelic.HasFiniteMapToP1
  \lean{AlgebraicGeometry.Adelic.exists_finiteMorphismToP1}
  % [papaioannou] Cor. 1.7-1.8 (a transcendental x has finitely many zeros/poles; [F:K(x)] finite); [cff] Thm. 12.6
  \source{papaioannou-algebraic-rr, cff-curves-function-fields}
  There is a nonconstant element \(x \in K \setminus k\) with
  \([K : k(x)] = n < \infty\); equivalently a finite \(k\)-morphism
  \(\pi \colon C \to \mathbb{P}^1_k\) of degree \(n\).
\end{proposition}
\begin{proof}
  As \(C\) is a geometrically integral curve over \(k\), its function
  field \(K/k\) has transcendence degree one. Choose any transcendental
  \(x \in K \setminus k\). Then \(k(x) \subseteq K\) is a subfield with
  \(K/k(x)\) algebraic; because \(K\) is a finitely generated field
  extension of \(k\) of transcendence degree one, \(K/k(x)\) is a finite
  extension, say of degree \(n = [K : k(x)]\)
  (Papaioannou Cor.~1.7--1.8: a transcendental element has finitely many
  zeros and poles and \(K/k(x)\) is finite). The inclusion
  \(k(x) = k(\mathbb{P}^1) \hookrightarrow K\) is the function-field map
  of a dominant, hence (properness of \(C\)) finite, morphism
  \(\pi \colon C \to \mathbb{P}^1_k\) of degree \(n\).
  \paragraph{Gate.} The existence of the finite map is packaged as the
  single-field class \(\mathrm{HasFiniteMapToP1}\, C\), carrying the
  witness \((x, hx : x \notin k, n, [K : k(x)] = n)\); a proved instance
  (transcendence-degree-one \(\Rightarrow\) such an \(x\)) is deferred,
  so Tier 1 finiteness verifies against the hypothesis.
\end{proof}

\begin{proposition}
[Charts are module-finite over the \(\mathbb{P}^1\) charts]
  \label{thm:adelic_chart_finite_over_P1}
  \lean{AlgebraicGeometry.Adelic.chart_finite_over_P1}
  \uses{thm:adelic_exists_finiteMorphismToP1, thm:adelic_chartRing_isDedekindDomain}
  % [papaioannou] sec. 1.3 (integrality of function fields); integral closure of k[x] in the finite extension K is module-finite
  \source{papaioannou-algebraic-rr}
  Take \(x\) as in \Cref{thm:adelic_exists_finiteMorphismToP1} and the
  cover \(U_0 = \pi^{-1}(\mathbb{A}^1_x)\),
  \(U_1 = \pi^{-1}(\mathbb{A}^1_{1/x})\). Then
  \(A_0\) is the integral closure of \(k[x]\) in \(K\), \(A_1\) that of
  \(k[1/x]\), and \(A_{01}\) that of \(k[x, 1/x]\); each is module-finite
  over \(k[x]\), \(k[1/x]\), \(k[x, 1/x]\) respectively, and free of rank
  \(n\) after localising to the generic point.
\end{proposition}
\begin{proof}
  The affine line \(\mathbb{A}^1_x = \Spec k[x]\) has coordinate ring
  \(k[x]\), and \(U_0 = \pi^{-1}(\mathbb{A}^1_x)\) is affine with
  coordinate ring \(A_0 = \Gamma(U_0, \struct{C})\), a Dedekind domain
  with \(\operatorname{Frac} A_0 = K\)
  (\Cref{thm:adelic_chartRing_isDedekindDomain}). Since \(A_0\) is
  integrally closed with fraction field \(K\) and contains \(k[x]\), and
  is integral over \(k[x]\) (finite fibres of \(\pi\)), it is precisely
  the integral closure of \(k[x]\) in \(K\). The integral closure of the
  Dedekind domain \(k[x]\) in the finite extension \(K/k(x)\) is a
  finitely generated \(k[x]\)-module (finiteness of the ring of integers
  of a function field over its rational subfield); as \(k[x]\) is a
  principal ideal domain and \(A_0\) is torsion-free of generic rank
  \(n\), \(A_0\) is \(k[x]\)-free of rank \(n\). The same argument at
  \(1/x\) gives \(A_1\) free of rank \(n\) over \(k[1/x]\), and localising
  either at \(x\) gives \(A_{01}\) module-finite over
  \(k[x, 1/x] = k[x]_x\).
\end{proof}

\begin{lemma}
[The \(\mathbb{P}^1\) base cokernel vanishes]
  \label{thm:adelic_cokernel_P1_zero}
  \lean{AlgebraicGeometry.Adelic.cokernel_P1_zero}
  % [papaioannou] p. 7-8 example (K(x) has genus 0); [cff] Example 14.5 (K(x) is of genus 0)
  \source{papaioannou-algebraic-rr, cff-curves-function-fields}
  The untwisted cokernel of the standard cover of \(\mathbb{P}^1_k\)
  vanishes:
  \[
    k[x, x^{-1}] \big/ \bigl(k[x] + k[x^{-1}]\bigr) \;=\; 0,
    \qquad\text{i.e.}\qquad H^1(\mathbb{P}^1_k, \struct{\mathbb{P}^1}) = 0.
  \]
\end{lemma}
\begin{proof}
  As a \(k\)-vector space \(k[x, x^{-1}]\) has basis the Laurent
  monomials \(\{x^m : m \in \mathbb{Z}\}\). Every monomial \(x^m\) with
  \(m \ge 0\) lies in \(k[x]\), and every monomial with \(m \le 0\) lies
  in \(k[x^{-1}]\); since each integer is \(\ge 0\) or \(\le 0\), the sum
  \(k[x] + k[x^{-1}]\) contains every basis monomial and hence equals
  \(k[x, x^{-1}]\). Therefore the quotient is zero. This is the genus-zero
  statement for \(\mathbb{P}^1_k\).
\end{proof}

\begin{theorem}
[KEYSTONE: the base cokernel is finite-dimensional]
  \label{thm:adelic_finiteDimensional_cokernel_zero}
  \lean{AlgebraicGeometry.Adelic.finiteDimensional_cokernel_zero}
  \uses{thm:adelic_chart_finite_over_P1, thm:adelic_cokernel_P1_zero, def:adelic_coverH1}
  % [papaioannou] Thm. 1.17-1.18 (finiteness of genus and of dim(A_F/(A_F(D)+F)))
  \source{papaioannou-algebraic-rr}
  The base cover cohomology is finite-dimensional:
  \[
    \mathrm{FiniteDimensional}\ k\ \bigl(A_{01}/(A_0 + A_1)\bigr)
    = \mathrm{FiniteDimensional}\ k\ \check{H}^1(0).
  \]
\end{theorem}
\begin{proof}
  Work with the \(\mathbb{P}^1\)-adapted cover of
  \Cref{thm:adelic_chart_finite_over_P1}: \(A_0\) is \(k[x]\)-free of
  rank \(n\), \(A_1\) is \(k[x^{-1}]\)-free of rank \(n\), and
  \(A_{01} = (A_0)_x\) is \(k[x, x^{-1}]\)-free of rank \(n\). Both
  \(A_0\) and \(A_1\) are \(k[x, x^{-1}]\)-lattices in the
  \(n\)-dimensional \(k(x)\)-vector space \(K\), so they become equal
  after inverting \(x\): \((A_0)_x = A_{01} = (A_1)_x\). By the structure
  theorem for a pair of lattices over the principal ideal domain
  \(k[x^{-1}]\) (Smith normal form / elementary divisors), there is a
  \(k(x)\)-basis \(e_1, \dots, e_n\) of \(K\) and integers
  \(m_1, \dots, m_n \in \mathbb{Z}\) with
  \[
    A_0 = \bigoplus_{j=1}^n k[x]\, e_j,
    \qquad
    A_1 = \bigoplus_{j=1}^n k[x^{-1}]\, x^{-m_j} e_j .
  \]
  Consequently the cokernel decomposes as a direct sum of rank-one
  \(\mathbb{P}^1\) cokernels:
  \[
    A_{01}/(A_0 + A_1)
    \;\cong\;
    \bigoplus_{j=1}^n
      k[x, x^{-1}] \big/ \bigl(k[x] + x^{-m_j} k[x^{-1}]\bigr).
  \]
  For \(m_j \le 0\) the summand is \(k[x, x^{-1}]/(k[x] + k[x^{-1}]) = 0\)
  by \Cref{thm:adelic_cokernel_P1_zero} (the submodule
  \(x^{-m_j}k[x^{-1}] \supseteq k[x^{-1}]\) when \(m_j \le 0\), so the sum
  is already all of \(k[x,x^{-1}]\)). For \(m_j \ge 1\) the submodule
  \(k[x] + x^{-m_j} k[x^{-1}]\) contains every monomial \(x^p\) with
  \(p \ge 0\) or \(p \le -m_j\), so the summand is spanned by the finitely
  many classes of \(x^{-1}, x^{-2}, \dots, x^{-(m_j - 1)}\), of dimension
  \(m_j - 1 < \infty\). Each summand is finite-dimensional, hence so is
  the finite direct sum, which is
  \(\check{H}^1(0)\). (The total dimension \(\sum_{j : m_j \ge 1}(m_j -1)\)
  is the genus of \(C\); we need only its finiteness here.)
\end{proof}

\begin{theorem}
[Tier 1 deliverable: \(\genus C\) is honest]
  \label{thm:adelic_module_finite_genus}
  \lean{AlgebraicGeometry.Adelic.module_finite_genus}
  \uses{thm:adelic_HModuleOne_linearEquiv_cokernel, thm:adelic_finiteDimensional_cokernel_zero, def:genus}
  The first cohomology of the structure sheaf is finite over \(k\):
  \[
    \Module.\Finite\ k\ \bigl(\HModule\, k\, (\toModuleKSheaf C)\, 1\bigr).
  \]
  Consequently \(\genus C = \finrank_k
  \bigl(\HModule\, k\, (\toModuleKSheaf C)\, 1\bigr)\) is a genuine
  natural number.
\end{theorem}
\begin{proof}
  The keystone \Cref{thm:adelic_finiteDimensional_cokernel_zero} gives
  \(\mathrm{FiniteDimensional}\ k\ \check{H}^1(0)\); the bridge
  \Cref{thm:adelic_HModuleOne_linearEquiv_cokernel} is a \(k\)-linear
  isomorphism
  \(\HModule\, k\, (\toModuleKSheaf C)\, 1 \simeq_k \check{H}^1(0)\).
  Finite-dimensionality transports across a linear isomorphism, so
  \(\HModule\, k\, (\toModuleKSheaf C)\, 1\) is finite over \(k\). As
  \(\genus C\) is defined (\Cref{def:genus}) as the \(\finrank\) of this
  module, its value is the honest dimension rather than the junk \(0\)
  that \(\finrank\) returns on a module not known to be finite.
\end{proof}

\section{Tier 2 --- the \(\chi\)-ledger}
\label{sec:RiemannRoch_Adelic_chi}

\begin{proposition}
[Twist exact sequence]
  \label{thm:adelic_twist_exact_sequence}
  \lean{AlgebraicGeometry.Adelic.twist_exact_sequence}
  \uses{def:adelic_coverH1, thm:adelic_dim_adele_quotient_eq_degree}
  % [papaioannou] Lemma 1.20 (0 -> L(D2)/L(D1) -> A(D2)/A(D1) -> (A(D2)+F)/(A(D1)+F) -> 0); [cff] Lemma 13.1(2)
  \source{papaioannou-algebraic-rr, cff-curves-function-fields}
  For Weil divisors \(D \le D'\) there is a four-term exact sequence of
  \(k\)-vector spaces
  \[
    0 \longrightarrow L(D')/L(D)
      \longrightarrow \mathcal{Q}(D, D')
      \longrightarrow \check{H}^1(D)
      \longrightarrow \check{H}^1(D')
      \longrightarrow 0,
  \]
  where \(\mathcal{Q}(D, D') = H^0\bigl(C, \struct{C}(D')/\struct{C}(D)\bigr)\)
  is the space of global sections of the skyscraper quotient sheaf.
\end{proposition}
\begin{proof}
  Because \(D \le D'\), the section spaces enlarge:
  \(\Gamma(U, \struct{C}(D)) \subseteq \Gamma(U, \struct{C}(D'))\) for
  each of \(U \in \{U_0, U_1, U_0 \cap U_1\}\). The two-term Čech
  complexes \(C^\bullet(D) \hookrightarrow C^\bullet(D')\) form a short
  exact sequence of complexes
  \(0 \to C^\bullet(D) \to C^\bullet(D') \to Q^\bullet \to 0\) with
  quotient complex
  \(Q^0 = C^0(D')/C^0(D)\), \(Q^1 = C^1(D')/C^1(D)\). The associated long
  exact cohomology sequence (snake lemma) is the six-term
  \[
    0 \to L(D) \to L(D') \to H^0(Q^\bullet)
      \to \check{H}^1(D) \to \check{H}^1(D') \to H^1(Q^\bullet) \to 0 .
  \]
  Now \(Q^\bullet\) is the Čech complex of the quotient sheaf
  \(\mathcal{T} = \struct{C}(D')/\struct{C}(D)\), a torsion sheaf
  supported on the finite set where \(D' \ne D\)
  (\Cref{thm:adelic_finite_support_order}). A skyscraper (finite-support
  torsion) sheaf on a curve is flasque, hence acyclic on every cover, so
  \(H^1(Q^\bullet) = 0\); equivalently the map \(Q^0 \to Q^1\) is
  surjective. Surjectivity is visible directly: an element
  \(\alpha \in C^1(D') = \Gamma(U_0 \cap U_1, \struct{C}(D'))\) has, at
  each of the finitely many points \(P \in U_0 \cap U_1\) where \(D'\)
  allows a larger pole than \(D\), a principal part that can be absorbed
  into whichever chart \(U_0\) or \(U_1\) contains \(P\) --- a section of
  \(\struct{C}(D')\) on that affine chart realising the prescribed pole
  (weak approximation on the affine chart) --- so
  \(\alpha \equiv f_0 - f_1 \pmod{C^1(D)}\) with
  \(f_i \in \Gamma(U_i, \struct{C}(D'))\). With \(H^1(Q^\bullet) = 0\) the
  six-term sequence truncates to the asserted four-term sequence, and
  \(H^0(Q^\bullet) = \ker(Q^0 \to Q^1) = H^0(C, \mathcal{T}) =
  \mathcal{Q}(D, D')\).
\end{proof}

\begin{proposition}
[Local computation: quotient dimension equals degree]
  \label{thm:adelic_dim_adele_quotient_eq_degree}
  \lean{AlgebraicGeometry.Adelic.dim_adele_quotient_eq_degree}
  \uses{def:adelic_pointValuation, thm:adelic_finite_support_order, def:divisor_degree}
  % [papaioannou] Lemma 1.19 (dim A(D2)/A(D1) = deg D2 - deg D1); [cff] Lemma 13.1(1) and (C1)-(C3); [scheidler] p. 24
  \source{papaioannou-algebraic-rr, cff-curves-function-fields, scheidler-function-fields}
  For \(D \le D'\) the skyscraper section space is finite-dimensional
  with
  \[
    \dim_k \mathcal{Q}(D, D')
      \;=\; \deg(D' - D)
      \;=\; \sum_P \bigl(D'(P) - D(P)\bigr)\,\deg P .
  \]
\end{proposition}
\begin{proof}
  The quotient sheaf \(\mathcal{T} = \struct{C}(D')/\struct{C}(D)\) is
  supported on the finite set \(\{P : D'(P) > D(P)\}\)
  (\Cref{thm:adelic_finite_support_order}), and its global sections split
  as the direct sum of stalks
  \(\mathcal{Q}(D, D') = \bigoplus_P \mathcal{T}_P\) (independent local
  conditions at distinct points). It therefore suffices to compute each
  local length, and by additivity we reduce to a single unit step
  \(D' = D + P_0\). Fix a uniformiser \(t\) at \(P_0\), i.e.\
  \(v_{P_0}(t) = 1\). The map
  \[
    \varphi \colon \Gamma(U, \struct{C}(D + P_0)) \longrightarrow \kappa(P_0),
    \qquad
    f \longmapsto \bigl(t^{\,D(P_0) + 1} f\bigr)(P_0),
  \]
  (the leading Laurent coefficient at \(P_0\)) is \(k\)-linear with kernel
  exactly \(\Gamma(U, \struct{C}(D))\): \(\varphi(f) = 0\) iff
  \(v_{P_0}(t^{D(P_0)+1} f) \ge 1\) iff \(v_{P_0}(f) \ge -D(P_0)\). It is
  surjective onto \(\kappa(P_0)\) because \(P_0 \in U\) affine, so the
  evaluation \(\Gamma(U, \struct{C}) \to \kappa(P_0)\) is surjective and
  multiplying representatives by \(t^{-(D(P_0)+1)}\) realises any residue
  class. Hence
  \(\mathcal{T}_{P_0} \cong \Gamma(U, \struct{C}(D + P_0))/
  \Gamma(U, \struct{C}(D)) \cong \kappa(P_0)\), of \(k\)-dimension
  \([\kappa(P_0) : k] = \deg P_0\)
  (\Cref{def:adelic_pointValuation}; cf.\ the graded-piece computation
  \(\dim_k \mathfrak p^{i}/\mathfrak p^{i+1} = \deg P\)). Summing the unit
  steps over the finite support gives
  \(\dim_k \mathcal{Q}(D, D') = \sum_P (D'(P) - D(P))\deg P = \deg(D' - D)\),
  the degree of \Cref{def:divisor_degree}.
\end{proof}

\begin{theorem}
[\(\chi\)-KEYSTONE: \(\chi(D) = \deg D + 1 - g\)]
  \label{thm:adelic_chi_eq_chi_zero_add_degree}
  \lean{AlgebraicGeometry.Adelic.chi_eq_chi_zero_add_degree}
  \uses{thm:adelic_twist_exact_sequence, thm:adelic_dim_adele_quotient_eq_degree, thm:adelic_module_finite_genus, def:genus, inst:Scheme_instIsHModuleHomFinite_toModuleKSheaf}
  % [papaioannou] Thm. 1.18 (dim D = deg D + 1 - g + i(D)); [cff] Cor. 12.11, Thm. 14.1
  \source{papaioannou-algebraic-rr, cff-curves-function-fields}
  The Euler characteristic \(\chi(D) = \ell(D) - i(D)\) is additive in the
  degree:
  \[
    \chi(D') - \chi(D) = \deg(D' - D)
    \quad\text{for } D \le D',
    \qquad\text{hence}\qquad
    \chi(D) = \deg D + 1 - g .
  \]
\end{theorem}
\begin{proof}
  By \Cref{thm:adelic_module_finite_genus} the base cohomology
  \(\check{H}^1(0)\) is finite-dimensional; the twist sequences below
  bound every \(\ell(D), i(D)\) by finite numbers, so all dimensions
  occurring are finite and the alternating-sum identity for an exact
  sequence of finite-dimensional spaces applies. Apply it to the
  four-term sequence of \Cref{thm:adelic_twist_exact_sequence} for
  \(D \le D'\):
  \[
    \dim\!\bigl(L(D')/L(D)\bigr) - \dim \mathcal{Q}(D, D')
      + \dim \check{H}^1(D) - \dim \check{H}^1(D') = 0 .
  \]
  With \(\dim(L(D')/L(D)) = \ell(D') - \ell(D)\),
  \(\dim \mathcal{Q}(D, D') = \deg(D' - D)\)
  (\Cref{thm:adelic_dim_adele_quotient_eq_degree}), and
  \(\dim \check{H}^1(D) = i(D)\), this rearranges to
  \[
    \bigl(\ell(D') - i(D')\bigr) - \bigl(\ell(D) - i(D)\bigr)
      = \deg(D' - D),
    \qquad\text{i.e.}\qquad \chi(D') - \chi(D) = \deg(D' - D).
  \]
  Every divisor \(D\) satisfies \(D - D_- \le D \le D_+\) with
  \(0 \le D_+\); applying the additivity to \(0 \le D_+\) and to
  \(D \le D + D_-\) and adding, one gets \(\chi(D) = \chi(0) + \deg D\)
  for all \(D\) (additivity of \(\chi - \deg\) as a function on the free
  group of divisors, which agrees at \(0\) and is additive on each unit
  step, hence is \(\chi(0)\) plus the homomorphism \(\deg\)). Finally
  \(\chi(0) = \ell(0) - i(0) = 1 - g\): the global sections
  \(L(0) = \Gamma(U_0, \struct{C}) \cap \Gamma(U_1, \struct{C}) =
  H^0(C, \struct{C})\) equal the constants \(k\), of dimension
  \(\ell(0) = 1\), because \(C\) is geometrically integral and proper
  (the finiteness/one-dimensionality of global sections is the input
  \Cref{inst:Scheme_instIsHModuleHomFinite_toModuleKSheaf}); and
  \(i(0) = \dim_k \check{H}^1(0) = g\) is the genus
  (\Cref{def:genus,thm:adelic_HModuleOne_linearEquiv_cokernel}). Hence
  \(\chi(D) = \deg D + 1 - g\).
\end{proof}

\begin{corollary}
[Riemann inequality; all \(\check{H}^1(D)\) finite]
  \label{thm:adelic_riemann_inequality}
  \lean{AlgebraicGeometry.Adelic.riemann_inequality}
  \uses{thm:adelic_chi_eq_chi_zero_add_degree, thm:adelic_module_finite_genus, thm:adelic_twist_exact_sequence, thm:adelic_dim_adele_quotient_eq_degree}
  % [papaioannou] Thm. 1.17 (Riemann inequality deg D - dim D <= g - 1); [cff] Thm. 12.10
  \source{papaioannou-algebraic-rr, cff-curves-function-fields}
  For every Weil divisor \(D\), the space \(\check{H}^1(D)\) is
  finite-dimensional and
  \[
    \deg D + 1 - g \;\le\; \ell(D) .
  \]
\end{corollary}
\begin{proof}
  Finiteness: for \(0 \le D''\) with \(D \le D''\) and \(0 \le D''\), the
  last map of the twist sequence
  (\Cref{thm:adelic_twist_exact_sequence}) makes
  \(\check{H}^1(D) \to \check{H}^1(D'')\) surjective with kernel a
  quotient of \(\mathcal{Q}(D, D'')\), which is finite
  (\Cref{thm:adelic_dim_adele_quotient_eq_degree}); and
  \(\check{H}^1(0) \to \check{H}^1(D'')\) is surjective, so
  \(\dim \check{H}^1(D'') \le \dim \check{H}^1(0) = g\). Hence
  \(\dim \check{H}^1(D) \le g + \deg(D'' - D) < \infty\). Inequality: from
  \(\chi(D) = \ell(D) - i(D) = \deg D + 1 - g\)
  (\Cref{thm:adelic_chi_eq_chi_zero_add_degree}) and
  \(i(D) = \dim_k \check{H}^1(D) \ge 0\) we obtain
  \(\ell(D) = \deg D + 1 - g + i(D) \ge \deg D + 1 - g\).
\end{proof}

\section{Tier 3 --- Serre duality and vanishing (gated)}
\label{sec:RiemannRoch_Adelic_duality}

This tier upgrades the Riemann inequality to the sharp Riemann--Roch
theorem via a residue pairing, and derives the curve-level vanishing
\(i(D) = 0\) for \(\deg D > 2g - 2\). Its two deepest nodes are gated by
single-field \(\mathrm{Prop}\)-classes; the earlier tiers do not depend on
them.

\begin{definition}
[Residue pairing on the cover cohomology]
  \label{thm:adelic_residuePairing}
  \notready  % gate class AlgebraicGeometry.Adelic.HasResiduePairing
  \lean{AlgebraicGeometry.Adelic.residuePairing}
  \uses{def:adelic_sectionOfDivisor, def:adelic_coverH1}
  % [vater] Def. 1.58 / Thm. 1.59 (local components omega_P; omega(alpha) = sum_P omega_P(alpha_P)); Thm. 5.66 (residues res_{P}); [papaioannou] Def. 1.13
  \source{vater-weil-differentials, papaioannou-algebraic-rr}
  Let \(\Omega_K\) be the space of Weil differentials of \(K/k\)
  (\(k\)-linear forms on the adeles vanishing on some
  \(\mathbb{A}_K(D) + K\); locally the Kähler differentials
  \(\Omega_{K/k}\), a one-dimensional \(K\)-vector space). There is a
  perfect \(k\)-bilinear residue pairing
  \[
    \langle\,\cdot\,,\,\cdot\,\rangle \colon
      \mathbb{A}_K \times \Omega_K \to k,
    \qquad
    (\alpha, \omega) \mapsto \sum_P \operatorname{Res}_P(\alpha_P\, \omega),
  \]
  where \(\operatorname{Res}_P\) is the local residue at \(P\). It
  satisfies the residue theorem \(\sum_P \operatorname{Res}_P(f\omega) = 0\)
  for \(f \in K\), and therefore descends to a pairing
  \[
    \check{H}^1(D) \times \Omega_K(K_C - D) \longrightarrow k,
  \]
  where \(\Omega_K(E) = \{\omega : \operatorname{div}\omega \ge E\}\) and
  \(K_C\) is a canonical divisor.
\end{definition}
\begin{proof}
  A Weil differential \(\omega\) has, at each place \(P\), a local
  component \(\omega_P \colon K \to k\) (Vater Def.~1.58) and evaluates on
  an adele by the finite sum
  \(\omega(\alpha) = \sum_P \omega_P(\alpha_P)\) (Vater Thm.~1.59), which
  is the residue sum \(\sum_P \operatorname{Res}_P(\alpha_P \omega)\)
  under the identification of local components with residues (Vater
  Thm.~5.66). Vanishing on principal adeles,
  \(\omega(f, f, \dots) = \sum_P \operatorname{Res}_P(f\omega) = 0\), is
  the definition of a Weil differential vanishing on \(K\), i.e.\ the
  residue theorem. If \(\alpha \in \mathbb{A}_K(D)\) and
  \(\operatorname{div}\omega \ge K_C - D\) then
  \(v_P(\alpha_P) + v_P(\omega) \ge -D(P) + (K_C - D)(P) \ge 0\) fails to
  contribute a pole, so \(\operatorname{Res}_P(\alpha_P\omega) = 0\) and
  the pairing kills \(\mathbb{A}_K(D) + K\), descending to
  \(\check{H}^1(D) \times \Omega_K(K_C - D) \to k\).
  \paragraph{Gate.} The Weil-differential space \(\Omega_K\), the local
  residues \(\operatorname{Res}_P\), and the residue theorem over an
  arbitrary \(k\) require Kähler-differential and trace machinery not yet
  assembled; they are threaded as the single-field class
  \(\mathrm{HasResiduePairing}\, C\) bundling the pairing and its residue
  theorem, discharged later by \(\texttt{inferInstance}\).
\end{proof}

\begin{theorem}
[Serre duality: \(i(D) = \ell(K_C - D)\)]
  \label{thm:adelic_serreDuality}
  \notready  % gate class AlgebraicGeometry.Adelic.HasSerreDuality
  \lean{AlgebraicGeometry.Adelic.serreDuality}
  \uses{thm:adelic_residuePairing, thm:adelic_chi_eq_chi_zero_add_degree}
  % [papaioannou] Thm. 1.23 (dim_F Omega_F = 1), Thm. 1.26 (mu: L(W-D) ~ Omega_F(D), i(D) = dim(W-D)); [cff] Thm. 13.5; [vater] Thm. 1.48, 1.55
  \source{papaioannou-algebraic-rr, cff-curves-function-fields, vater-weil-differentials}
  For a canonical divisor \(K_C = \operatorname{div}\omega\) of a nonzero
  Weil differential \(\omega\), the map \(x \mapsto x\omega\) induces a
  \(k\)-linear isomorphism
  \[
    \mu \colon L(K_C - D) \;\xrightarrow{\ \sim\ }\; \Omega_K(K_C - D)
    \;\simeq_k\; \check{H}^1(D)^\ast,
  \]
  so that \(i(D) = \dim_k \check{H}^1(D) = \ell(K_C - D)\).
\end{theorem}
\begin{proof}
  The residue pairing (\Cref{thm:adelic_residuePairing}) identifies
  \(\Omega_K(K_C - D)\) with the annihilator of
  \(\mathbb{A}_K(D) + K\), i.e.\ with
  \((\mathbb{A}_K/(\mathbb{A}_K(D) + K))^\ast = \check{H}^1(D)^\ast\)
  (Papaioannou Lemma 1.22; Vater Thm.~1.48), so
  \(\dim_k \Omega_K(K_C - D) = i(D)\). For the first isomorphism, since
  \(\dim_K \Omega_K = 1\) (Papaioannou Thm.~1.23; Vater Thm.~1.48) every
  differential is \(x\omega\) for a unique \(x \in K\); and
  \(\operatorname{div}(x\omega) = \operatorname{div} x + K_C\) (Papaioannou
  Thm.~1.25; Vater Lemma 1.53), so \(x\omega\) lies in
  \(\Omega_K(K_C - D)\) iff \(\operatorname{div} x + K_C \ge K_C - D\) iff
  \(\operatorname{div} x + D \ge 0\) iff \(x \in L(D)\); replacing \(D\) by
  \(K_C - D\) gives that \(\mu\) restricts to a bijection
  \(L(K_C - D) \to \Omega_K(K_C - D)\) (Papaioannou Thm.~1.26; cff
  Thm.~13.5). Composing, \(i(D) = \ell(K_C - D)\).
  \paragraph{Gate.} The construction of \(\Omega_K\), the
  one-dimensionality \(\dim_K \Omega_K = 1\), and the non-degeneracy of
  the pairing are threaded as the single-field class
  \(\mathrm{HasSerreDuality}\, C\) bundling the equivalence \(\mu\);
  Tier 3 consumers verify against it and it is closed by
  \(\texttt{inferInstance}\) once \(\mathrm{HasResiduePairing}\) is discharged.
\end{proof}

\begin{corollary}
[Vanishing: \(i(D) = 0\) for \(\deg D > 2g - 2\)]
  \label{thm:adelic_h1_vanishing}
  \lean{AlgebraicGeometry.Adelic.h1_vanishing}
  \uses{thm:adelic_serreDuality, thm:adelic_chi_eq_chi_zero_add_degree}
  % [papaioannou] Cor. 1.28 (deg W = 2g-2), Cor. 1.29 (RR sharp for deg D >= 2g-1); [cff] Cor. 14.2, 14.3
  \source{papaioannou-algebraic-rr, cff-curves-function-fields}
  If \(\deg D > 2g - 2\) then \(\check{H}^1(\struct{C}(D))\) is a
  subsingleton: \(i(D) = 0\). This is the curve-level replacement for
  Castelnuovo--Mumford boundedness in the Quot endgame.
\end{corollary}
\begin{proof}
  A canonical divisor has degree \(\deg K_C = 2g - 2\)
  (from Riemann--Roch at \(D = 0\) and \(D = K_C\): Papaioannou Cor.~1.28;
  cff Cor.~14.2). By Serre duality \Cref{thm:adelic_serreDuality},
  \(i(D) = \ell(K_C - D)\). When \(\deg D > 2g - 2\) we have
  \(\deg(K_C - D) = (2g - 2) - \deg D < 0\). A divisor of negative degree
  has \(\ell = 0\): if \(\ell(E) > 0\) then some \(f \in L(E)\) makes
  \(\operatorname{div} f + E \ge 0\) an effective divisor of degree
  \(\deg E\) (principal divisors have degree \(0\),
  \Cref{thm:adelic_finite_support_order}), forcing \(\deg E \ge 0\).
  Hence \(\ell(K_C - D) = 0\), so \(i(D) = 0\) and
  \(\check{H}^1(\struct{C}(D))\) is a subsingleton.
\end{proof}

\begin{theorem}
[Riemann--Roch]
  \label{thm:adelic_riemannRoch}
  \lean{AlgebraicGeometry.Adelic.riemannRoch}
  \uses{thm:adelic_chi_eq_chi_zero_add_degree, thm:adelic_serreDuality}
  % [papaioannou] Thm. 1.27 (dim D = deg D + 1 - g + dim(W - D)); [cff] Thm. 14.1; [vater] Thm. 1.55
  \source{papaioannou-algebraic-rr, cff-curves-function-fields, vater-weil-differentials}
  For every Weil divisor \(D\) and canonical divisor \(K_C\),
  \[
    \ell(D) - \ell(K_C - D) \;=\; \deg D + 1 - g .
  \]
\end{theorem}
\begin{proof}
  Immediate from the two keystones: the \(\chi\)-ledger
  \Cref{thm:adelic_chi_eq_chi_zero_add_degree} gives
  \(\ell(D) - i(D) = \deg D + 1 - g\), and Serre duality
  \Cref{thm:adelic_serreDuality} gives \(i(D) = \ell(K_C - D)\);
  substituting the second into the first yields the identity.
\end{proof}
